\section{收益的边界：什么时候不该分离}
\label{sec:boundary}

文献普遍报告正收益，但这里有一个发布偏差：失败的部署不会写成论文。
故本节的证据需要主动找——找那些\textbf{专门问「何时不划算」}的工作。

\subsection{一条直接的判据}

最清晰的表述来自对 Prefill--Decode--Attention--FFN 四段分离的模拟研究：
更细的分离粒度「不保证值得它的复杂度」，收益取决于模型架构，
也取决于负载的 prefill/output 比率\texttt{[c-boundary]}。

这句话的价值在于它把问题从「分离好不好」改成了
「在当前架构与负载比下，分离好不好」。前者无解，后者可测。

\subsection{四类使分离失效的条件}

\noindent\textbf{① 传输成本超过收益。} 这是最直接的一类。由第
\ref{sec:transfer} 节的数据，同一份 250\,MB 的 KV 经 RDMA 传输需 10\,ms。
若 decode 的单步时延目标是 50\,ms 量级，传输就吃掉了 20\%——
而分离带来的配置优化能否补回这 20\%，要看收益侧有多大。

\noindent\textbf{② 两阶段资源画像接近。} 小模型、短上下文、单卡可容纳
全部权重与 KV 时，prefill 与 decode 的最优配置差别很小。此时分离只增加
传输与复杂度，收益趋近于零。

\noindent\textbf{③ 负载本身很轻。} 算力闲置时，配置是否最优无关紧要——
反正都够快。分离的收益本质上来自「在资源紧张时把资源用得更对」，
资源不紧张时这个收益无从兑现。

\noindent\textbf{④ 上下文极短导致固定开销主导。} 这一条最反直觉：
KV 绝对量小时，按字节数算传输应该很快，但建连、内存注册、同步这些
固定开销不随消息大小缩放。此时带宽不是瓶颈，固定开销才是。

\subsection{一个明确的反例}

必须单独列出这条，因为它直接挑战「分离更优」的默认假设。

有工作报告：相对 PD 分离与 PD 聚合，某方案的 TTFT 最多降低 13.2×，
TPOT 最多降低 1.69×\texttt{[c-boundary]}。也就是说，
\textbf{存在一种方案，在 TTFT 上比分离好 13 倍}。
如果只读「分离 vs 不分离」的对比，会完全看不到这个量级的改进空间——
因为分离本身可能不是这个维度上的最优解。

这条证据的意义不在于否定分离，而在于提醒：
「分离 vs 混部」是一个二选一的问题框架，而真实的设计空间里
还有第三、第四种做法（如按注意力与 FFN 再拆、如用 CXL 换载体）。
把问题设成二选一，会系统性地漏掉这些。

\subsection{判据汇总}

\begin{table}[t]
\centering
\caption{分离的决策判据。四项全满足才建议分离；任一项不满足时，
应优先考虑 chunked prefill 等零传输成本的方案。}
\label{tab:criteria}
\small
\begin{tabular}{@{}p{0.35cm}p{3.0cm}p{2.5cm}@{}}
\toprule
\# & 条件 & 检查方式 \\
\midrule
1 & 两阶段最优配置显著不同 & 分别做单阶段压测，比最优并行度与 batch \\
2 & 存在高带宽 KV 通路 & 同机 NVLink 或跨机 RDMA，而非共享存储 \\
3 & 负载足够重 & 资源利用率已高，配置优化能兑现 \\
4 & KV 规模足够大 & 传输固定开销占比可忽略 \\
\bottomrule
\end{tabular}
\end{table}

这四项不是充分条件，但是必要条件——任一项不满足时，
分离的收益很难为正。它们也是可操作的：每一项都能在部署前测出来。
